\section{Phản ứng hóa học và biến đổi chất}

\subsection{Biến đổi vật lí và biến đổi hóa học}

Trong tự nhiên và trong phòng thí nghiệm, chất có thể thay đổi về trạng thái, hình dạng, kích thước hoặc có thể biến đổi thành chất khác. Muốn phân biệt hai trường hợp này, cần xét xem \textbf{có tạo thành chất mới hay không}.

\begin{ghinho}{Phân biệt hai loại biến đổi}
\textbf{Biến đổi vật lí} là biến đổi không tạo thành chất mới. Chất ban đầu vẫn được giữ nguyên về thành phần hóa học, chỉ thay đổi trạng thái, hình dạng, kích thước hoặc mức độ phân chia.

\textbf{Biến đổi hóa học} là biến đổi có tạo thành chất mới. Khi đó đã xảy ra phản ứng hóa học.
\end{ghinho}

Ví dụ: nước đá tan thành nước lỏng là biến đổi vật lí; iron để lâu ngoài không khí ẩm tạo gỉ là biến đổi hóa học. Khi đánh giá một hiện tượng, không nên chỉ dựa vào việc ``có thay đổi'' mà phải hỏi: \emph{sau quá trình, có chất mới được tạo thành hay không?}

\begin{vidu}
Xét các hiện tượng sau: (1) nghiền nhỏ đá vôi; (2) nung đá vôi tạo vôi sống và khí carbon dioxide; (3) hòa tan đường vào nước; (4) đốt cháy khí methane. Hãy phân loại các hiện tượng trên thành biến đổi vật lí và biến đổi hóa học.
\end{vidu}
\begin{loigiai}
Các hiện tượng (1) và (3) không tạo chất mới nên là biến đổi vật lí.

Ở hiện tượng (2), calcium carbonate tạo thành calcium oxide và carbon dioxide:
\[
\ce{CaCO3 ->[t^\circ] CaO + CO2}.
\]
Ở hiện tượng (4), methane phản ứng với oxygen tạo carbon dioxide và nước:
\[
\ce{CH4 + 2O2 -> CO2 + 2H2O}.
\]
Hai hiện tượng này tạo chất mới nên là biến đổi hóa học.

\textbf{Vậy,} (1), (3) là biến đổi vật lí; (2), (4) là biến đổi hóa học.
\end{loigiai}

\subsection{Phản ứng hóa học, chất đầu và sản phẩm}

Phản ứng hóa học là quá trình biến đổi từ chất này thành chất khác. Các chất có mặt trước phản ứng và bị biến đổi được gọi là \textbf{chất đầu} (hay chất phản ứng); các chất mới tạo thành được gọi là \textbf{sản phẩm}.

Một phương trình hóa học cho biết chất nào tham gia, chất nào tạo thành và tỉ lệ số mol giữa các chất. Chẳng hạn:
\[
\ce{Zn + 2HCl -> ZnCl2 + H2}.
\]
Trong phản ứng trên, Zn và \ce{HCl} là chất đầu; \ce{ZnCl2} và \ce{H2} là sản phẩm.

\begin{ghinho}{Dấu hiệu nhận biết phản ứng hóa học}
Một số dấu hiệu thường gặp khi phản ứng xảy ra:
\begin{itemize}
\item xuất hiện chất khí;
\item xuất hiện hoặc biến mất kết tủa;
\item màu sắc thay đổi;
\item nhiệt độ thay đổi, phát sáng hoặc có ngọn lửa;
\item xuất hiện mùi mới.
\end{itemize}
Các dấu hiệu trên là căn cứ thực nghiệm quan trọng nhưng cần kết hợp với kiến thức về chất. Ví dụ, nước sôi tạo bọt khí nhưng chỉ là biến đổi trạng thái, không phải phản ứng hóa học.
\end{ghinho}

\begin{vidu}
Nhỏ dung dịch hydrochloric acid vào một mẩu calcium carbonate. Quan sát thấy chất rắn tan dần và có nhiều bọt khí thoát ra. Viết phương trình hóa học và nêu dấu hiệu chứng tỏ có phản ứng xảy ra.
\end{vidu}
\begin{loigiai}
Phương trình hóa học:
\[
\ce{CaCO3 + 2HCl -> CaCl2 + CO2 + H2O}.
\]
Dấu hiệu rõ nhất là có khí \ce{CO2} thoát ra và chất rắn \ce{CaCO3} tan dần.

\textbf{Vậy,} đã xảy ra phản ứng giữa \ce{CaCO3} và \ce{HCl}, tạo \ce{CaCl2}, \ce{CO2} và \ce{H2O}.
\end{loigiai}

\subsection{Năng lượng trong phản ứng hóa học}

Phản ứng hóa học thường kèm theo sự trao đổi năng lượng với môi trường. Nếu phản ứng truyền năng lượng ra môi trường dưới dạng nhiệt thì gọi là \textbf{phản ứng tỏa nhiệt}; nếu phản ứng nhận nhiệt từ môi trường thì gọi là \textbf{phản ứng thu nhiệt}.

\begin{ghinho}{Nhận biết phản ứng tỏa nhiệt và thu nhiệt}
\begin{itemize}
\item \textbf{Tỏa nhiệt:} môi trường xung quanh thường nóng lên. Ví dụ: đốt cháy than, xăng, methane; phản ứng của quicklime với nước.
\item \textbf{Thu nhiệt:} phản ứng cần được cung cấp nhiệt liên tục hoặc làm môi trường lạnh đi. Ví dụ: nhiệt phân calcium carbonate.
\end{itemize}
Không nên suy luận chỉ từ việc ``phải châm lửa lúc đầu''. Nhiều phản ứng tỏa nhiệt vẫn cần năng lượng ban đầu để khởi động, sau đó phản ứng tự tiếp tục và giải phóng năng lượng.
\end{ghinho}

\begin{vidu}
Khi cho nước vào quicklime, khối chất rắn nóng lên rõ rệt:
\[
\ce{CaO + H2O -> Ca(OH)2}.
\]
Phản ứng trên thuộc loại tỏa nhiệt hay thu nhiệt? Giải thích.
\end{vidu}
\begin{loigiai}
Sau khi trộn, nhiệt độ của hệ tăng, chứng tỏ phản ứng đã truyền năng lượng ra môi trường dưới dạng nhiệt.

\textbf{Vậy,} phản ứng \ce{CaO + H2O -> Ca(OH)2} là phản ứng tỏa nhiệt.
\end{loigiai}

\subsection{Định luật bảo toàn khối lượng}

Trong phản ứng hóa học, các nguyên tử không tự sinh ra và cũng không tự mất đi; chúng chỉ sắp xếp lại để tạo thành các chất mới. Vì vậy, trong một hệ kín, tổng khối lượng các chất trước phản ứng bằng tổng khối lượng các chất sau phản ứng.

\begin{ghinho}{Định luật bảo toàn khối lượng}
Với phản ứng tổng quát
\[
\ce{A + B -> C + D},
\]
trong một hệ kín:
\[
m_{\ce{A}}+m_{\ce{B}}=m_{\ce{C}}+m_{\ce{D}}.
\]
Nếu phản ứng sinh khí và tiến hành trong bình hở, số chỉ của cân có thể giảm vì khí thoát ra khỏi hệ. Điều này không làm định luật bảo toàn khối lượng mất đúng; khi đó phần khí thoát ra đã không còn nằm trên cân.
\end{ghinho}

\begin{vidu}
Cho \SI{10.0}{\gram} calcium carbonate phản ứng hoàn toàn với \SI{14.6}{\gram} hydrochloric acid trong một hệ kín. Sau phản ứng thu được \SI{11.1}{\gram} calcium chloride và nước cùng khí carbon dioxide. Biết khối lượng nước tạo thành là \SI{1.8}{\gram}. Tính khối lượng \ce{CO2}.
\end{vidu}
\begin{loigiai}
Theo định luật bảo toàn khối lượng:
\[
10.0+14.6=11.1+1.8+m_{\ce{CO2}}.
\]
Suy ra
\[
m_{\ce{CO2}}=11.7\ \text{g}.
\]

\textbf{Vậy,} khối lượng khí carbon dioxide tạo thành là \SI{11.7}{\gram}.
\end{loigiai}

\begin{tuluyen}
Phân loại các hiện tượng sau thành biến đổi vật lí hoặc biến đổi hóa học: nước hoa bay hơi; iron nóng chảy; copper bị phủ lớp màu xanh ngoài không khí ẩm; đốt cháy ethanol; hòa tan sodium chloride vào nước.
\dapso{Vật lí: nước hoa bay hơi, iron nóng chảy, hòa tan sodium chloride; hóa học: copper bị phủ lớp màu xanh, đốt cháy ethanol.}
\begin{loigiaiBT}
Nước hoa bay hơi và iron nóng chảy chỉ thay đổi trạng thái; hòa tan sodium chloride không tạo chất mới. Copper để lâu ngoài không khí ẩm tạo các hợp chất mới trên bề mặt; ethanol cháy tạo carbon dioxide và nước.

\textbf{Vậy,} ba hiện tượng đầu nêu trên thuộc nhóm vật lí như đã phân loại; hai hiện tượng tạo chất mới thuộc biến đổi hóa học.
\end{loigiaiBT}
\end{tuluyen}

\begin{tuluyen}
Một học sinh cân một cốc chứa dung dịch acid và một viên đá vôi, được tổng khối lượng \SI{152.4}{\gram}. Sau khi thả viên đá vôi vào cốc và chờ phản ứng kết thúc, khối lượng hệ trên cân còn \SI{150.2}{\gram}. Giải thích vì sao số chỉ của cân giảm và tính khối lượng khí đã thoát ra.
\dapso{Khí sinh ra thoát khỏi cốc; khối lượng khí thoát ra là \SI{2.2}{\gram}.}
\begin{loigiaiBT}
Cốc là hệ hở nên khí tạo thành thoát khỏi phần vật chất đang đặt trên cân. Khối lượng khí thoát ra:
\[
152.4-150.2=\SI{2.2}{\gram}.
\]
\textbf{Vậy,} khối lượng khí đã thoát ra là \SI{2.2}{\gram}; định luật bảo toàn khối lượng vẫn đúng nếu tính cả lượng khí này.
\end{loigiaiBT}
\end{tuluyen}
